\section{Related Work}
\label{sec:related}

\paragraph{Multimodal meme understanding.}
Memes convey meaning through the interplay of an image and overlaid text. The Hateful Memes
challenge~\cite{kiela2020hateful} showed that unimodal models fail when either modality alone is benign. Memotion
(SemEval-2020 Task~8)~\cite{sharma2020memotion} introduced the humour, sarcasm, offensiveness, motivation and
overall-sentiment taxonomy that CMBan adopts. Most systems pair a pretrained text encoder with a pretrained vision
encoder, often CLIP~\cite{radford2021clip}, and fuse the two representations.

\paragraph{Bangla meme resources.}
Bangla memes have recently received dedicated datasets.
\begin{itemize}
  \item MemoSen~\cite{hossain2022memosen} annotates 4{,}417 memes for three-way sentiment and reports that
  multimodal models outperform unimodal ones.
  \item MUTE~\cite{hossain2022mute} targets hateful memes in Bangla and code-mixed text.
  \item BHM~\cite{hossain2024bhm} adds target identification for hateful memes.
  \item BanglaAbuseMeme~\cite{das2023banglaabusememe} annotates 4{,}043 memes for abuse, vulgarity, sarcasm,
  sentiment and target.
  \item MIMOSA~\cite{ahsan2024mimosa} annotates aggression targets.
  \item CMBan~\cite{cmban2025} is the first resource built around \emph{cartoon} memes. It annotates 2{,}641
  Bangla and code-mixed memes along the five Memotion dimensions and identifies more than 800 recurring cartoon
  templates.
\end{itemize}
All of these works evaluate models on held-out \emph{clean} data. None asks whether the predictions survive the
spelling, script and image variations that memes undergo as they spread.

\paragraph{Robustness evaluation.}
Controlled corruption benchmarks such as ImageNet-C~\cite{hendrycks2019benchmarking} established severity-graded
perturbations as a standard way to measure visual robustness. In NLP, behavioural testing (CheckList;
\cite{ribeiro2020checklist}) probes models with targeted input transformations. Code-switching benchmarks such as
GLUECoS~\cite{khanuja2020gluecos} show that multilingual encoders degrade on mixed-language input. For Bangla, the
dominant real-world variation is \emph{Romanised} (``Banglish'') writing, with highly inconsistent spelling. CLIP is
known to read text rendered in images, which makes it vulnerable to typographic manipulation~\cite{goh2021multimodal}.
That property is directly relevant to memes, where the text \emph{is} part of the image.

\paragraph{Explainability for multimodal models.}
Gradient-based attribution methods such as Integrated Gradients~\cite{sundararajan2017ig} and
Grad-CAM~\cite{selvaraju2017gradcam,chattopadhay2018gradcampp} are widely used to inspect which tokens and image
regions drive a decision. Prior Bangla meme work has used explanations mainly to justify predictions. We instead use
attribution \emph{comparatively}: we measure how modality reliance and visual attention shift between the clean and
perturbed versions of the same meme.

\paragraph{Positioning.}
ToonSense is, to our knowledge, the first study of the \emph{robustness} of multimodal models on Bangla cartoon
memes. It combines a multi-encoder benchmark, Bangla-specific linguistic perturbations (colloquial transliteration,
Romanised spelling variation, code-mixing) with visual corruptions, before/after attribution analysis, and a source
attribution of errors.
